\begin{examplebox}

	\textbf{\textcolor{CExample}{例 1（基础）}}

	\textbf{\textcolor{CExample}{题目：}}
	在正项等比数列 $\{a_n\}$ 中，已知 $a_2 = 2$，$a_4 = 8$，求 $a_3$。

	\textbf{\textcolor{CExample}{解题步骤：}}
	\begin{description}[style=nextline, leftmargin=2.8em, labelsep=0.5em, itemsep=0.45em]
		\item[\textbf{第一步（识别条件）：}] 由 $a_2$、$a_3$、$a_4$ 是连续三项，故 $a_3$ 是 $a_2$ 与 $a_4$ 的等比中项，有 $a_3^2 = a_2 a_4$。
			\par\textbf{理由：} 代入 $n=3$ 时的等比中项公式。

		\item[\textbf{第二步（代入计算）：}] $a_3^2 = 2 \times 8$，得 $a_3^2 = 16$，故 $a_3 = \pm 4$。
			\par\textbf{理由：} 方程 $x^2=16$ 有两个实数根。

		\item[\textbf{第三步（取舍符号）：}] 因为数列是正项等比数列，所有项均为正数，故 $a_3 = 4$。
			\par\textbf{理由：} 正项数列排除负值。
	\end{description}

	\textbf{\textcolor{CExample}{关键结论：}}
	$\boxed{a_3 = 4}$。

	\medskip
	\textbf{\textcolor{CExample}{例 2（稍变形）}}

	\textbf{\textcolor{CExample}{题目：}}
	已知三个数 $2$、$x$、$8$ 成等比数列，求 $x$ 的值。

	\textbf{\textcolor{CExample}{解题步骤：}}
	\begin{description}[style=nextline, leftmargin=2.8em, labelsep=0.5em, itemsep=0.45em]
		\item[\textbf{第一步（套用公式）：}] 由等比中项性质得 $x^2 = 2 \times 8$，即 $x^2 = 16$。
			\par\textbf{理由：} 若三数成等比，则中间平方等于两端积。

		\item[\textbf{第二步（解方程）：}] $x = \pm 4$。
			\par\textbf{理由：} 对 $16$ 开平方得 $\pm 4$。

		\item[\textbf{第三步（确定答案）：}] 题中无额外符号条件，故 $x = 4$ 或 $x = -4$ 均正确。
			\par\textbf{理由：} 等比中项可取正负。
	\end{description}

	\textbf{\textcolor{CExample}{关键结论：}}
	$\boxed{x = \pm 4}$。

	\medskip
	\textbf{\textcolor{CExample}{例 3（综合应用）}}

	\textbf{\textcolor{CExample}{题目：}}
	在等比数列 $\{a_n\}$ 中，$a_1 = 3$，$a_3 = 12$，且 $a_2$ 与 $a_3$ 同号，求公比 $q$。

	\textbf{\textcolor{CExample}{解题步骤：}}
	\begin{description}[style=nextline, leftmargin=2.8em, labelsep=0.5em, itemsep=0.45em]
		\item[\textbf{第一步（求中间项）：}]
			\[
				\begin{aligned}
					a_2^2 &= a_1 a_3 \\
					&= 3 \times 12 \\
					&= 36
			\end{aligned}
			\]
			故 $a_2 = \pm 6$。
			\par\textbf{理由：} 等比中项公式。

		\item[\textbf{第二步（符号判定）：}] $a_3 = 12 > 0$，且 $a_2$ 与 $a_3$ 同号，则 $a_2 > 0$，所以 $a_2 = 6$。
			\par\textbf{理由：} 同号取正。

		\item[\textbf{第三步（求公比）：}]
			\[
				\begin{aligned}
					q &= \dfrac{a_2}{a_1} \\
					&= \dfrac{6}{3} \\
					&= 2.
			\end{aligned}
			\]
			\par\textbf{理由：} 公比等于相邻两项之比。
	\end{description}

	\textbf{\textcolor{CExample}{关键结论：}}
	$\boxed{q = 2}$。

\end{examplebox}
